\documentclass[a4paper,12pt]{article}

\usepackage{fontspec}
\setmainfont{STIX Two Text}
\usepackage[height=1.6\baselineskip]{mahjong}
\usepackage[skip=10pt plus2pt]{parskip}
\usepackage{csquotes}

\setlength{\parindent}{0pt}

\newcommand{\mj}[1]{\mahjong{#1}}
%\newcommand{\mmj}[1]{\parbox[0][3\baselineskip]{\linewidth}{\mahjong{#1}}}
\newcommand{\mmj}[1]{\begin{displayquote}\mahjong{#1}\end{displayquote}}
\newcommand{\nl}{\newline}

\begin{document}
\section{Han Scorers}

These are features of a hand that give Han, one of the measures 
that is used to calculate the score a hand provides.

\section{Non-Yaku}

These are features of a hand that provide Han, but do not 
count as Yaku and themselves cannot make a hand valid.

\subsection{Dora: 1 Han}

At the start of the game one of the tiles on the Dead Wall 
is revealed as a Dora Indicator. Since the Dead Wall is 
7 tiles long, this would look like e.g.
\mmj{XX4XXXXp}
Whichever tile comes after the Dora Indicator is a Dora tile, 
so in the above example \mj{5p} would be a Dora tile.
The sequence loops so if \mj{9m} is the Dora Indicator 
then \mj{1m} is a Dora tile.

The Honor tiles also have an order, the Wind tiles are in 
the order of play \mj{1234z}. Similarly the Dragon tiles 
have an order \mj{567z}, but you can remember the sequence 
alphabetically, green to red to white.

\subsection{Kandora: 1 Han}

After a player calls Kan, the next tile of the Dead Wall 
is flipped and becomes another Dora indicator. Continuing 
from the example above the wall might now show 
\mmj{XX4p7sXXX} 
Which means that \mj{5p} and \mj{8s} are both Dora tiles.

\subsection{Akadora: 1 Han}

In Riichi Mahjong there are often a set of red 5s. Some sets 
have one in each suit, and some have one in each except two in 
the circle suit. These tiles look like \mj{0m0s0p}.

\subsection{Stacking Dora}

One tile can be a Dora tile multiple times, and thus be 
worth multiple Han. For example, if the Dora Indicators showed 
\mmj{22XXXs}
then \mj{3s} would be worth 2 Han. Similarly if the 
Dora Indicator was
\mmj{4XXXXm}
then \mj{0m} would be worth 2 Han.

\section{Yaku}

Yaku are features of a hand that provide Han. Importantly 
a hand needs to have a Yaku in order to be able to win in 
Riichi Mahjong. 

\subsection{Riichi: 1 Han}

A closed hand in Tenpai can call Riichi, typically the 
first discard made after calling Riichi is turned sideways. 
Once in Riichi you cannot change your hand, only drawing 
a tile and discarding it if it doesn't complete the hand.

\subsection{Double Riichi: 2 Han}

A Riichi called on the first turn is a Double Riichi.

\subsection{Ippatsu: 1 Han}

If after calling Riichi you win before your next discard you 
have this Yaku. Any player making a call (Chii, Pon, Kan) 
invalidates this Yaku.

\subsection{Yakuhai: 1 Han}

Having a set of any three dragons, your seat wind or the 
prevalent wind counts as the Yakuhai Yaku. An example with the 
red dragons might look like:
\mmj{234p7p222m789m-77'7z-7p}
And if you were sitting in the south seat the following hand 
would be Yakuhai:
\mmj{45p99p123s333m-2'22z-6p}

\end{document}
